% Front matter: About this edition; then book PDF pp. 2, 9--14 (About the
% Author, Preface for Students, Preface for Instructors, Acknowledgments).
% The license and attribution are taken from the title page (PDF p. 1).

\chapter*{About this Edition}
\addcontentsline{toc}{chapter}{About this Edition}

This is a re-typeset edition of \emph{Linear Algebra Done Right}, fourth
edition, by Sheldon Axler, \textcopyright\ 2024 Sheldon Axler. It follows the
open access version of the book dated 13 July 2026. The print version of the
book is published by Springer.

The book is licensed under the terms of the Creative Commons
Attribution-NonCommercial 4.0 International License (CC BY-NC 4.0,
\url{https://creativecommons.org/licenses/by-nc/4.0}), which permits any
noncommercial use, sharing, adaptation, distribution and reproduction in any
medium or format, as long as appropriate credit is given to the original author
and source, a link to the license is provided, and changes are indicated. This
edition is made for noncommercial use only.

This edition changes typography and layout only. The text, the mathematics,
the exercises, and the numbering of definitions, results, and examples are
Axler's. The photographs that open the chapters and their captions are not
reproduced. The indexes (Symbol Index and Index) and the page references of the
original are omitted because the pagination differs. The grouping of the
chapters into four parts is an addition of this edition. Errors introduced in the
re-typesetting are not the author's.

The solutions to the exercises of Chapters 1--4 and Sections 5A--5B are by
blanketism (\url{https://github.com/blanketism/Axler-Solutions}); they are not
part of Axler's book and appear in the Solutions chapter.

% ---------------------------------------------------------------- PDF p. 2
\chapter*{About the Author}
\addcontentsline{toc}{chapter}{About the Author}

Sheldon Axler received his undergraduate degree from Princeton University,
followed by a PhD in mathematics from the University of California at
Berkeley.

As a postdoctoral Moore Instructor at MIT, Axler received a university-wide
teaching award. He was then an assistant professor, associate professor, and
professor at Michigan State University, where he received the first
J.~Sutherland Frame Teaching Award and the Distinguished Faculty Award.

Axler received the Lester R. Ford Award for expository writing from the
Mathematical Association of America in 1996, for a paper that eventually
expanded into this book. In addition to publishing numerous research papers, he
is the author of six mathematics textbooks, ranging from freshman to graduate
level. Previous editions of this book have been adopted as a textbook at over
375 universities and colleges and have been translated into three languages.

Axler has served as Editor-in-Chief of the \emph{Mathematical Intelligencer}
and Associate Editor of the \emph{American Mathematical Monthly}. He has been a
member of the Council of the American Mathematical Society and of the Board of
Trustees of the Mathematical Sciences Research Institute. He has also served on
the editorial board of Springer's series Undergraduate Texts in Mathematics,
Graduate Texts in Mathematics, Universitext, and Springer Monographs in
Mathematics.

Axler is a Fellow of the American Mathematical Society and has been a
recipient of numerous grants from the National Science Foundation.

Axler joined San Francisco State University as chair of the Mathematics
Department in 1997. He served as dean of the College of Science \& Engineering
from 2002 to 2015, when he returned to a regular faculty appointment as a
professor in the Mathematics Department.

\begin{figure}[!htbp]\centering
\includegraphics[page=2,trim=23.7 86.7 33.05 461,clip,width=\linewidth]{source-book.pdf}
\par\smallskip{\itshape The author and his cat Moon.}\par
{\footnotesize Carrie Heeter, Bishnu Sarangi}
\end{figure}

\emph{Cover equation}: Formula for the $n^{\text{th}}$ Fibonacci number.
Exercise 21 in Section 5D uses linear algebra to derive this formula.
\[ F_n = 1/\sqrt5\biggl[\biggl(\frac{1+\sqrt5}{2}\biggr)^{\!n}
         -\biggl(\frac{1-\sqrt5}{2}\biggr)^{\!n}\biggr] \]

% ---------------------------------------------------------------- PDF p. 9
\chapter*{Preface for Students}
\addcontentsline{toc}{chapter}{Preface for Students}

You are probably about to begin your second exposure to linear algebra. Unlike
your first brush with the subject, which probably emphasized Euclidean spaces
and matrices, this encounter will focus on abstract vector spaces and linear
maps. These terms will be defined later, so don't worry if you do not know what
they mean. This book starts from the beginning of the subject, assuming no
knowledge of linear algebra. The key point is that you are about to immerse
yourself in serious mathematics, with an emphasis on attaining a deep
understanding of the definitions, theorems, and proofs.

You cannot read mathematics the way you read a novel. If you zip through a
page in less than an hour, you are probably going too fast. When you encounter
the phrase ``as you should verify'', you should indeed do the verification,
which will usually require some writing on your part. When steps are left out,
you need to supply the missing pieces. You should ponder and internalize each
definition. For each theorem, you should seek examples to show why each
hypothesis is necessary. Discussions with other students should help.

As a visual aid, definitions and notations are in yellow boxes (and have a
title beginning with the word ``\textsf{definition}'' or ``\textsf{notation}'')
and theorems are in blue boxes (in color versions of the book). Each theorem
has an informal descriptive name.

Please check the website below for additional information about the book,
including a link to videos that are freely available to accompany the book.

Your suggestions, comments, and corrections are most welcome.

Best wishes for success and enjoyment in learning linear algebra!

\medskip
\noindent\begin{tabular}{@{}l@{}}
Sheldon Axler\\
San Francisco State University\\[3pt]
website: \url{https://linear.axler.net}\\
e-mail: \href{mailto:linear@axler.net}{linear@axler.net}
\end{tabular}

% ------------------------------------------------------------ PDF pp. 10--13
\chapter*{Preface for Instructors}
\addcontentsline{toc}{chapter}{Preface for Instructors}

You are about to teach a course that will probably give students their second
exposure to linear algebra. During their first brush with the subject, your
students probably worked with Euclidean spaces and matrices. In contrast, this
course will emphasize abstract vector spaces and linear maps.

The title of this book deserves an explanation. Most linear algebra textbooks
use determinants to prove that every linear operator on a finite-dimensional
complex vector space has an eigenvalue. Determinants are difficult,
nonintuitive, and often defined without motivation. To prove the theorem about
existence of eigenvalues on complex vector spaces, most books must define
determinants, prove that a linear operator is not invertible if and only if its
determinant equals $0$, and then define the characteristic polynomial. This
tortuous (torturous?) path gives students little feeling for why eigenvalues
exist.

In contrast, the simple determinant-free proofs presented here (for example,
see \ref{ax:5.19}) offer more insight. Once determinants have been moved to the
end of the book, a new route opens to the main goal of linear
algebra---understanding the structure of linear operators.

This book starts at the beginning of the subject, with no prerequisites other
than the usual demand for suitable mathematical maturity. A few examples and
exercises involve calculus concepts such as continuity, differentiation, and
integration. You can easily skip those examples and exercises if your students
have not had calculus. If your students have had calculus, then those examples
and exercises can enrich their experience by showing connections between
different parts of mathematics.

Even if your students have already seen some of the material in the first few
chapters, they may be unaccustomed to working exercises of the type presented
here, most of which require an understanding of proofs.

Here is a chapter-by-chapter summary of the highlights of the book:
\begin{itemize}
\item Chapter 1: Vector spaces are defined in this chapter, and their basic
  properties are developed.
\item Chapter 2: Linear independence, span, basis, and dimension are defined in
  this chapter, which presents the basic theory of finite-dimensional vector
  spaces.
\item Chapter 3: This chapter introduces linear maps. The key result here is
  the fundamental theorem of linear maps: if $T$ is a linear map on $V$, then
  $\dim V=\dim\nul T+\dim\range T$. Quotient spaces and duality are topics in
  this chapter at a higher level of abstraction than most of the book; these
  topics can be skipped (except that duality is needed for tensor products in
  Section 9D).
\item Chapter 4: The part of the theory of polynomials that will be needed to
  understand linear operators is presented in this chapter. This chapter
  contains no linear algebra. It can be covered quickly, especially if your
  students are already familiar with these results.
\item Chapter 5: The idea of studying a linear operator by restricting it to
  small subspaces leads to eigenvectors in the early part of this chapter. The
  highlight of this chapter is a simple proof that on complex vector spaces,
  eigenvalues always exist. This result is then used to show that each linear
  operator on a complex vector space has an upper-triangular matrix with
  respect to some basis. The minimal polynomial plays an important role here
  and later in the book. For example, this chapter gives a characterization of
  the diagonalizable operators in terms of the minimal polynomial. Section 5E
  can be skipped if you want to save some time.
\item Chapter 6: Inner product spaces are defined in this chapter, and their
  basic properties are developed along with tools such as orthonormal bases and
  the Gram--Schmidt procedure. This chapter also shows how orthogonal
  projections can be used to solve certain minimization problems. The
  pseudoinverse is then introduced as a useful tool when the inverse does not
  exist. The material on the pseudoinverse can be skipped if you want to save
  some time.
\item Chapter 7: The spectral theorem, which characterizes the linear operators
  for which there exists an orthonormal basis consisting of eigenvectors, is
  one of the highlights of this book. The work in earlier chapters pays off
  here with especially simple proofs. This chapter also deals with positive
  operators, isometries, unitary operators, matrix factorizations, and
  especially the singular value decomposition, which leads to the polar
  decomposition and norms of linear maps.
\item Chapter 8: This chapter shows that for each operator on a complex vector
  space, there is a basis of the vector space consisting of generalized
  eigenvectors of the operator. Then the generalized eigenspace decomposition
  describes a linear operator on a complex vector space. The multiplicity of an
  eigenvalue is defined as the dimension of the corresponding generalized
  eigenspace. These tools are used to prove that every invertible linear
  operator on a complex vector space has a square root. Then the chapter gives
  a proof that every linear operator on a complex vector space can be put into
  Jordan form. The chapter concludes with an investigation of the trace of
  operators.
\item Chapter 9: This chapter begins by looking at bilinear forms and showing
  that the vector space of bilinear forms is the direct sum of the subspaces of
  symmetric bilinear forms and alternating bilinear forms. Then quadratic forms
  are diagonalized. Moving to multilinear forms, the chapter shows that the
  subspace of alternating $n$-linear forms on an $n$-dimensional vector space
  has dimension one. This result leads to a clean basis-free definition of the
  determinant of an operator. For complex vector spaces, the determinant turns
  out to equal the product of the eigenvalues, with each eigenvalue included in
  the product as many times as its multiplicity. The chapter concludes with an
  introduction to tensor products.
\end{itemize}

This book usually develops linear algebra simultaneously for real and complex
vector spaces by letting $\F$ denote either the real or the complex numbers. If
you and your students prefer to think of $\F$ as an arbitrary field, then see
the comments at the end of Section 1A. I prefer avoiding arbitrary fields at
this level because they introduce extra abstraction without leading to any new
linear algebra. Also, students are more comfortable thinking of polynomials as
functions instead of the more formal objects needed for polynomials with
coefficients in finite fields. Finally, even if the beginning part of the
theory were developed with arbitrary fields, inner product spaces would push
consideration back to just real and complex vector spaces.

You probably cannot cover everything in this book in one semester. Going
through all the material in the first seven or eight chapters during a
one-semester course may require a rapid pace. If you must reach Chapter 9, then
consider skipping the material on quotient spaces in Section 3E, skipping
Section 3F on duality (unless you intend to cover tensor products in
Section 9D), covering Chapter 4 on polynomials in a half hour, skipping
Section 5E on commuting operators, and skipping the subsection in Section 6C on
the pseudoinverse.

A goal more important than teaching any particular theorem is to develop in
students the ability to understand and manipulate the objects of linear
algebra. Mathematics can be learned only by doing. Fortunately, linear algebra
has many good homework exercises. When teaching this course, during each class
I usually assign as homework several of the exercises, due the next class.
Going over the homework might take up significant time in a typical class.

Some of the exercises are intended to lead curious students into important
topics beyond what might usually be included in a basic second course in
linear algebra.

\subsection*{The author's top ten}

Listed below are the author's ten favorite results in the book, in order of
their appearance in the book. Students who leave your course with a good
understanding of these crucial results will have an excellent foundation in
linear algebra.
\begin{itemize}
\item any two bases of a vector space have the same length (\ref{ax:2.34})
\item fundamental theorem of linear maps (\ref{ax:3.21})
\item existence of eigenvalues if $\F=\C$ (\ref{ax:5.19})
\item upper-triangular form always exists if $\F=\C$ (\ref{ax:5.47})
\item Cauchy--Schwarz inequality (\ref{ax:6.14})
\item Gram--Schmidt procedure (\ref{ax:6.32})
\item spectral theorem (\ref{ax:7.29} and \ref{ax:7.31})
\item singular value decomposition (\ref{ax:7.70})
\item generalized eigenspace decomposition theorem when $\F=\C$
  (\ref{ax:8.22})
\item dimension of alternating $n$-linear forms on $V$ is $1$ if $\dim V=n$
  (\ref{ax:9.37})
\end{itemize}

\subsection*{Major improvements and additions for the fourth edition}

\begin{itemize}
\item Over 250 new exercises and over 70 new examples.
\item Increasing use of the minimal polynomial to provide cleaner proofs of
  multiple results, including necessary and sufficient conditions for an
  operator to have an upper-triangular matrix with respect to some basis (see
  Section 5C), necessary and sufficient conditions for diagonalizability (see
  Section 5D), and the real spectral theorem (see Section 7B).
\item New section on commuting operators (see Section 5E).
\item New subsection on pseudoinverse (see Section 6C).
\item New subsections on QR factorization/Cholesky factorization (see
  Section 7D).
\item Singular value decomposition now done for linear maps from an inner
  product space to another (possibly different) inner product space, rather
  than only dealing with linear operators from an inner product space to itself
  (see Section 7E).
\item Polar decomposition now proved from singular value decomposition, rather
  than in the opposite order; this has led to cleaner proofs of both the
  singular value decomposition (see Section 7E) and the polar decomposition
  (see Section 7F).
\item New subsection on norms of linear maps on finite-dimensional inner
  product spaces, using the singular value decomposition to avoid even
  mentioning supremum in the definition of the norm of a linear map (see
  Section 7F).
\item New subsection on approximation by linear maps with lower-dimensional
  range (see Section 7F).
\item New elementary proof of the important result that if $T$ is an operator
  on a finite-dimensional complex vector space $V$, then there exists a basis
  of $V$ consisting of generalized eigenvectors of $T$ (see \ref{ax:8.9}).
\item New Chapter 9 on multilinear algebra, including bilinear forms, quadratic
  forms, multilinear forms, and tensor products. Determinants now are defined
  using a basis-free approach via alternating multilinear forms.
\item New formatting to improve the student-friendly appearance of the book.
  For example, the definition and result boxes now have rounded corners
  instead of right-angle corners, for a gentler look. The main font size has
  been reduced from 11 point to 10.5 point.
\end{itemize}

Please check the website below for additional links and information about the
book. Your suggestions, comments, and corrections are most welcome.

Best wishes for teaching a successful linear algebra class!

\marginnote{Contact the author, or Springer if the author is not available, for
permission for translations or other commercial reuse of the contents of this
book.}%
\medskip
\noindent\begin{tabular}{@{}l@{}}
Sheldon Axler\\
San Francisco State University\\[3pt]
website: \url{https://linear.axler.net}\\
e-mail: \href{mailto:linear@axler.net}{linear@axler.net}
\end{tabular}

% ---------------------------------------------------------------- PDF p. 14
\chapter*{Acknowledgments}
\addcontentsline{toc}{chapter}{Acknowledgments}

I owe a huge intellectual debt to all the mathematicians who created linear
algebra over the past two centuries. The results in this book belong to the
common heritage of mathematics. A special case of a theorem may first have been
proved long ago, then sharpened and improved by many mathematicians in
different time periods. Bestowing proper credit on all contributors would be a
difficult task that I have not undertaken. In no case should the reader assume
that any result presented here represents my original contribution.

Many people helped make this a better book. The three previous editions of
this book were used as a textbook at over 375 universities and colleges around
the world. I received thousands of suggestions and comments from faculty and
students who used the book. Many of those suggestions led to improvements in
this edition. The manuscript for this fourth edition was class tested at 30
universities. I am extremely grateful for the useful feedback that I received
from faculty and students during this class testing.

The long list of people who should be thanked for their suggestions would fill
up many pages. Lists are boring to read. Thus to represent all contributors to
this edition, I will mention only Noel Hinton, a graduate student at Australian
National University, who sent me more suggestions and corrections for this
fourth edition than anyone else. To everyone who contributed suggestions, let
me say how truly grateful I am to all of you. Many many thanks!

I thank Springer for providing me with help when I needed it and for allowing
me the freedom to make the final decisions about the content and appearance of
this book. Special thanks to the two terrific mathematics editors at Springer
who worked with me on this project---Loretta Bartolini during the first half of
my work on the fourth edition, and Elizabeth Loew during the second half of my
work on the fourth edition. I am deeply indebted to David Kramer, who did a
magnificent job of copyediting and prevented me from making many mistakes.

Extra special thanks to my fantastic partner Carrie Heeter. Her understanding
and encouragement enabled me to work intensely on this new edition. Our
wonderful cat Moon, whose picture appears on the \emph{About the Author} page,
provided sweet breaks throughout the writing process. Moon died suddenly due to
a blood clot as this book was being finished. We are grateful for five precious
years with him.

\begin{flushright}
Sheldon Axler
\end{flushright}
